\subsection{Identificación del problema técnico en la empresa}

\subsubsection{Problemática local}

\textbf{J\&P Periféricos S.A.C.}, como empresa proveedora de soluciones informáticas, ha encargado el desarrollo de un sistema de planillas que proyecta incorporar a su portafolio y ofrecer mediante suscripción a clientes empresariales. En este contexto, el problema técnico abordado consiste en desarrollar y validar una solución que integre el procesamiento de nóminas con la administración de la información laboral necesaria para su elaboración.
Según la entrevista al gerente general, las necesidades consideradas incluyen el registro de horas y licencias, la actualización de la información del personal y la integración tecnológica de los procesos. Estas observaciones orientan el levantamiento de requisitos de la solución. El proyecto no presupone que todas las empresas de la cartera presenten las mismas dificultades ni atribuye a J\&P incidentes internos de pago que no hayan sido documentados.
La propuesta comprende la centralización de registros personales, contractuales y remunerativos; la validación de entradas; el cálculo de nóminas; y la revisión de resultados mediante responsabilidades diferenciadas de Recursos Humanos y Contabilidad. Asimismo, incorpora control de acceso, historial de modificaciones y flujos automatizados para preparar archivos de pago y distribuir boletas. Su aporte se evaluará mediante pruebas funcionales y casos de cálculo definidos, distinguiendo los resultados comprobados de los beneficios esperados de su futura utilización comercial.

\subsubsection{Problemática nacional}

En el ámbito peruano, diversas organizaciones enfrentan serias complicaciones en la administración de su capital humano y nóminas debido a la persistente dependencia de hojas de cálculo convencionales y métodos manuales de registro. De acuerdo con \textcite{elperuano2026}, las inconsistencias más recurrentes en la gestión de compensaciones comprenden transferencias fuera de cronograma, dispersiones salariales incompletas y un cómputo inexacto de horas laborales, permisos y beneficios de ley, todo ello derivado de un bajo nivel de automatización en los flujos administrativos. Esta falta de estandarización técnica no solo expone a las organizaciones a contingencias administrativas y sanciones laborales, sino que afecta directamente la motivación, confianza y estabilidad de los colaboradores.

\subsubsection{Problemática internacional}

A escala global, la transición de procesos manuales hacia plataformas digitales integradas constituye un factor clave para la sostenibilidad y competitividad de las organizaciones. Según reportes del \textcite{bancomundial2024}, la adopción de herramientas tecnológicas incrementa de manera significativa los niveles de productividad y facilita el cumplimiento de normativas fiscales y laborales. No obstante, el organismo evidencia la existencia de una pronunciada brecha digital: mientras las corporaciones de gran envergadura aceleran la automatización de sus operaciones críticas, las pequeñas y medianas empresas avanzan con lentitud debido a limitaciones presupuestarias e infraestructura técnica limitada.
En el contexto regional, este desfase operativo se mantiene constante. Conforme señala \textcite{valenteneto2025} en un análisis para Ernst \& Young, la modernización y automatización de flujos operativos representa una de las mayores prioridades corporativas para afrontar la presión competitiva y responder con precisión a las demandas dinámicas del mercado. Sin embargo, las organizaciones de menor dimensión continúan atadas a tareas manuales y descentralizadas, lo cual eleva exponencialmente su margen de error en tareas de cálculo financiero e impacta de forma negativa en su sostenibilidad y competitividad.

\subsubsection{Formulación del Problema}

\noindent\textbf{Problema General:}
\begin{itemize}[leftmargin=*, label={}]
    \item ¿De qué manera la dispersión de la información laboral y las limitaciones de acceso, 
    control de modificaciones y trazabilidad afectan la confiabilidad del proceso de elaboración y 
    revisión de nóminas en las empresas destinatarias de la solución de J\&P Periféricos S.A.C.?
\end{itemize}

\vspace{0.3cm}
\noindent\textbf{Problemas Específicos:}
\begin{itemize}[
    leftmargin=*, label={\rule[0.2ex]{4.5pt}{4.5pt}}, topsep=4pt, itemsep=6pt ]
    \item ¿Cómo dificulta la dispersión de registros personales,
    contractuales y de incidencias la consolidación de información
    para elaborar la nómina?

    \item ¿Qué dificultades presenta el procedimiento actual
    para verificar los datos y cálculos antes de aprobar
    la planilla?

    \item ¿Qué limitaciones existen para restringir el acceso
    a la información salarial e identificar las modificaciones
    realizadas y sus responsables?

    \item ¿Qué dificultades presenta la preparación de archivos
    de pago y la distribución de boletas para controlar
    las versiones utilizadas y el resultado de las entregas?
\end{itemize}